\section{Full policy sweeps in the nonlinear investment economy}\label{sec:study54}
The directed certificate is now evaluated beyond an exact-null identity and beyond the final decision. This section uses the original continuous-law two-state investment economy, with state-dependent capacity, nonlinear costs and transitions, and acquired-state implementation. It does not substitute a finite-state model for that economy. The construction and direct-cost records were executed in the separately frozen R53 primary and adaptive studies; the present revision integrates them into the paper and replays their full deposited outputs. That replay adds neither new policies nor new independent samples.

\subsection{Design and the economic endpoint}
The primary design has horizons $T=2,3$, price parameter $p=1$, and two generators: the compiled witness construction and tensor FVI. Each service starts from the primitives and executes every rung of the original prefix up to the declared checkpoint. It checks that its reconstructed original policy has the same identity as the historical checkpoint. A five-bit observation contract in each coordinate then gives $32^2$ closed cells per date, with downward action spacing $1/4096$. Both methods use the same acquisition and robust-feasibility rules.

Starting from that complete acquired incumbent, each service executes exactly $T$ all-date passes. At every pass all new dates are computed against the same old complete incumbent and then adopted together. Each cell has nine robustly feasible proposed actions and eight interval covers of the full continuous-action range. Sixteen dyadic common-innovation bins enclose every unresolved expectation; the last innovation is integrated analytically. Rectangle extrema enclose discontinuous future actors. These calculations implement the directed certificate of Section~\ref{sec:directed54} with the verification gauge $h=0$. They are not a claim that a zero critic generates an effective policy, and their policy-evaluation cost is recorded.

The primary numerical endpoint is the actual expected cost of each complete returned policy under the original uniform initial law and continuous innovations. Conditional on the frozen policies, 65,536 common-path rows per horizon produce outward path-cost enclosures using forty-bit continuous-bin cells. The bounded empirical-Bernstein account applies to the endpoint samples under the declared independent-bin model. Dependence between policies through common paths is retained. The primary protocol allocates a simultaneous family error of $1/100$; the adaptive extension has a separate $1/100$ family. Combining the two families requires the corresponding union bound, rather than counting their repeated policies as independent economic objects.

\begin{table}[htbp]
\centering
\caption{Actual expected policy costs before and after full sweeps}\label{tab:primary-cost54}
{\small\setlength{\tabcolsep}{4pt}
\begin{tabular}{rrccc}
\hline
$T$ & Method & Initial cost & Final cost & Initial minus final \\
\hline
2 & W & [0.60258, 0.62351] & [0.59092, 0.61168] & [0.00665, 0.01684] \\
2 & F & [0.58999, 0.61078] & [0.58984, 0.61063] & [-0.00463, 0.00494] \\
3 & W & [0.71853, 0.74116] & [0.70898, 0.73147] & [0.00344, 0.01579] \\
3 & F & [0.70576, 0.72833] & [0.70576, 0.72833] & [-0.00596, 0.00597] \\
\hline
\end{tabular}}
\par\smallskip
{\footnotesize W is the compiled witness construction; F is tensor FVI. Every cost is for a complete acquired policy under the original continuous initial and innovation laws. Endpoints are rounded outward for display. Paired gain intervals, not subtraction of the displayed marginal intervals, determine gain signs. The policies and observations were frozen and executed in the R53 primary study.}
\end{table}

\begin{table}[htbp]
\centering
\caption{Paired witness-minus-FVI cost after every pass}\label{tab:cross-method54}
{\small\setlength{\tabcolsep}{4pt}
\begin{tabular}{rrc}
\hline
$T$ & Pass & Expected-cost difference \\
\hline
2 & 0 & [0.00757, 0.01775] \\
2 & 1 & [0.00151, 0.01148] \\
2 & 2 & [-0.00376, 0.00590] \\
3 & 0 & [0.00663, 0.01896] \\
3 & 1 & [0.00285, 0.01513] \\
3 & 2 & [-0.00105, 0.01112] \\
3 & 3 & [-0.00288, 0.00924] \\
\hline
\end{tabular}}
\par\smallskip
{\footnotesize A positive interval identifies higher witness cost. An interval containing zero is unresolved, not evidence of equality or witness superiority. The same paths couple policies within a horizon; no independence between these estimands is assumed.}
\end{table}


The distinction between cumulative gain and comparative ranking remains important. In the three-date experiment, the witness policy's gain after three full sweeps lies in $[0.00344390,0.01578829]$, with the displayed endpoints rounded outward. Its final witness-minus-FVI interval is $[-0.00287355,0.00923410]$. Thus the cumulative own-incumbent gain is identified, while the final cross-method difference is unresolved. The latter is not evidence of equality and not an identified reversal of the original ranking. Individual step gains can be statistically unresolved even though the deterministic certificate guarantees nonworsening expected cost at every true state.

For an installed witness controller with a replacement charge $\tau\ge0$, let $[l,u]$ be the corresponding cumulative gain interval. The net benefit of replacement is in $[l-\tau,u-\tau]$ on the same simultaneous event. Hence every $0\le\tau<l$ is supported as strictly beneficial when $l>0$, and every $\tau>u$ is ruled out as beneficial. The intermediate region is unresolved. This gives a full adoption-fee account for an actual improved policy, rather than only a reduced certificate for an unchanged policy. These are normalized model units, not an estimated installation charge or a calibrated welfare exercise.

\subsection{All-date acceptance and the original Bellman gap}
The improvement account records each proposal, the exact retained incumbent, the interval endpoints, the selected action, the continuous-action lower bound and the induced directed gap. Table~\ref{tab:sweeps54} gives every primary pass, including nonterminal changes. The supplementary tables report every date's positive contrast allowance and directed gap, the number of blocked and admitted candidates, and the comparison with scalar-width and symmetric-contrast gates. No failed or blocked cell is omitted.

\begin{table}[htbp]
\centering
\caption{All-date policy changes and certified gap recursion}\label{tab:sweeps54}
{\small\setlength{\tabcolsep}{4pt}
\begin{tabular}{rrr r r r r}
\hline
$T$ & Method & Pass & Changed & Nonterminal & Gap bound & Prefix (s) \\
\hline
2 & W & 1 & 590 & 17 & 3.72986 & 3.462 \\
2 & W & 2 & 265 & 265 & 0.08149 & 6.608 \\
2 & F & 1 & 653 & 1 & 3.59553 & 3.477 \\
2 & F & 2 & 0 & 0 & 0.06317 & 6.639 \\
3 & W & 1 & 612 & 44 & 5.11709 & 51.769 \\
3 & W & 2 & 282 & 282 & 3.65096 & 101.038 \\
3 & W & 3 & 135 & 135 & 0.22556 & 149.583 \\
3 & F & 1 & 64 & 0 & 4.87775 & 50.846 \\
3 & F & 2 & 0 & 0 & 3.40006 & 99.948 \\
3 & F & 3 & 0 & 0 & 0.09350 & 149.173 \\
\hline
\end{tabular}}
\par\smallskip
{\footnotesize Changed counts refer to distinct deployed cell-date actions in that pass, not to candidate comparisons. The gap is the directed all-state upper bound against the original continuous-action Bellman optimum. It need not decrease monotonically. Prefix work includes primitive construction and previous passes through their recorded array outputs; the final service and inference records are charged separately in the catalogue account.}
\end{table}


There are two distinct guarantees. A nonpositive directed upper advantage gives safety for an accepted action. The lower interval cover of all feasible actions gives the additional accuracy quantity $\varepsilon_t^k=\sup_C(U_t^k-L_t^k)$. The computed recursion is then $G_t^{k+1}=\beta_tG_{t+1}^k+\varepsilon_t^k$, initialized by a valid support bound. Its terms include cell uncertainty, actor ambiguity, action covering, integration and arithmetic through their actual endpoint propagation. They are not separately estimated from residual samples and then presumed negligible.

A coarse continuous-action cover or a wide future-actor range can make $G$ conservative. The theorem guarantees that the returned policy does not worsen, not that a particular numerical upper bound must decrease at every pass. Executing $T$ passes at nonzero tolerances also does not assert exact optimality: the remaining gap is the displayed sum of propagated directed errors. The exact $T$-pass optimality statement and its sharp coefficient have their own rational finite-model regression tests.

\subsection{A genuinely nonuniform conventional comparison}
The extension adds surplus-driven FVI at the same two-date primitive checkpoint. Its adaptation history and the number of nonuniform date models are recorded, together with every pilot and every earlier construction rung. It is not the earlier curvature rule that left its grid unchanged. Compiled witness and tensor FVI are reconstructed on the same extension runner, receive the same two complete passes, and reproduce their primary policy identities. A fresh path stream compares all three methods and all passes within a separately frozen statistical family.

\begin{table}[htbp]
\centering
\caption{Fresh actual-cost comparison with a nonuniform conventional method}\label{tab:adaptive54}
{\small\setlength{\tabcolsep}{4pt}
\begin{tabular}{rrcc}
\hline
Method & Pass & Actual expected cost & Gain from initial policy \\
\hline
W & 0 & [0.60337, 0.62425] & -- \\
W & 1 & [0.59714, 0.61789] & [0.00129, 0.01130] \\
W & 2 & [0.59173, 0.61245] & [0.00662, 0.01681] \\
F & 0 & [0.59081, 0.61156] & -- \\
F & 1 & [0.59065, 0.61141] & [-0.00462, 0.00494] \\
F & 2 & [0.59065, 0.61141] & [-0.00462, 0.00494] \\
A & 0 & [0.59068, 0.61143] & -- \\
A & 1 & [0.59060, 0.61135] & [-0.00470, 0.00486] \\
A & 2 & [0.59060, 0.61135] & [-0.00470, 0.00486] \\
\hline
\end{tabular}}
\par\smallskip
{\footnotesize A is surplus-driven FVI; W and F denote the same generators as above. All three are reconstructed in the same separately frozen cohort and receive two full improvement passes. W and F policy identities reproduce the primary cohort, while the cost paths are a fresh family. Nonuniformity and pilot work are recorded explicitly.}
\end{table}


This comparison tests an actual changed conventional representation, rather than merely a new label for the tensor baseline. Its adaptive grid nevertheless remains a tensor grid with nonuniform coordinates. It is therefore not presented as a sparse-grid or non-tensor high-dimensional demonstration. The earlier failed curvature rule, the adverse graded refinement, and the original direct-cost comparisons remain in their original tables and records.

\subsection{Complete catalogue work at common actual-cost targets}
A finite catalogue of returned policies supplies a precise work account. For each method we select the pass with the lowest simultaneous actual-cost upper endpoint. We charge that method's entire observed construction-and-improvement service, including all executed passes and durable records, plus the entire shared cost-inference service. Charging the full shared service to each eligible method is conservative and avoids assigning an unobserved marginal inference time to a method evaluated only jointly. The separate publication audit reports the additional replay and document-production overhead.

\begin{table}[htbp]
\centering
\caption{Complete-catalogue work to a verified actual-cost upper bound}\label{tab:catalogue54}
{\small\setlength{\tabcolsep}{4pt}
\begin{tabular}{rrr r r r r}
\hline
Cohort & Method & Pass & Cost upper & Service (s) & Inference (s) & Sum (s) \\
\hline
P2 & W & 2 & 0.61168 & 6.612 & 1.850 & 8.463 \\
P2 & F & 1 & 0.61063 & 6.643 & 1.850 & 8.493 \\
P3 & W & 3 & 0.73147 & 149.587 & 3.182 & 152.769 \\
P3 & F & 1 & 0.72833 & 149.178 & 3.182 & 152.359 \\
A2 & W & 2 & 0.61245 & 6.646 & 2.667 & 9.312 \\
A2 & F & 1 & 0.61141 & 6.579 & 2.667 & 9.245 \\
A2 & A & 1 & 0.61135 & 18.521 & 2.667 & 21.188 \\
\hline
\end{tabular}}
\par\smallskip
{\footnotesize For each method the selected pass minimizes its actual-cost upper endpoint. Work conservatively charges every own executed construction and sweep, plus the entire shared inference service. It is not a first-crossing or minimum-work sequential stopping estimate. Replay and document-production overhead are separately timed in the release audit. The exact finite threshold partition, eligible methods and least-recorded-work choices are deposited in the machine-readable audit.}
\end{table}


For a target $u$, the eligible methods are exactly those whose selected actual-cost upper bound is at most $u$. Taking the least recorded complete-catalogue work among eligible methods defines the observed finite frontier. Its breakpoints and method choices are deposited without rounding; the printed table is descriptive. This is a frontier for releasing a policy from an already executed finite catalogue. It is not a claim about the minimum work of an unexecuted sequential stopping algorithm, an optimized resource allocation, or a universal fastest generator. In particular, improving the witness incumbent does not by itself establish a better work-to-cost frontier than FVI.

The work tables separately identify construction time, pair-node counts, coupling work, actor queries, actor-table bytes, process peak memory, and the cumulative clocks. A common-innovation certificate avoids a stored table over all state pairs, but the work bound in Proposition~\ref{prop:complexity54} makes its horizon dependence explicit. These records turn a constructive mathematical certificate into an executed and costed verification method in the stated economy. They do not remove the tensor observation cover or establish a representation-specific scaling advantage.

\subsection{Controlled approximate-null and false-null tests}
The extension varies three departures from the exact action-null example: the ReLU direction, the deterministic action exposure, and a known action-dependent innovation-mean coefficient. Each departure ranges over $0$, $1/4096$, $1/256$, and $1/16$; the amplitudes are $1$, $16$, and $4096$. Two original incumbents yield 384 controlled cases, each evaluated on all $32^2$ closed cells. A verified positive allowance replaces the zero certificate whenever the altered coefficient is nonzero.

\begin{table}[htbp]
\centering
\caption{Controlled deviations from exact action-null structure}\label{tab:misspecification54}
{\small\setlength{\tabcolsep}{4pt}
\begin{tabular}{lr}
\hline
Quantity & Recorded value \\
\hline
Controlled cases & 384 \\
Closed cells per case & 1024 \\
Cell--case combinations & 393216 \\
Positive-bound cases & 378 \\
Safe changed cell--cases & 170374 \\
Harmful false-null cell--cases & 116185 \\
Maximum contrast allowance & 315.00000 \\
New statistical observations & 0 \\
\hline
\end{tabular}}
\par\smallskip
{\footnotesize The catalogue changes known ReLU directions, action exposures and action-dependent innovation means. Harm means a strictly positive lower true-cost difference for an action selected by the deliberately false-null rule. Corrected decisions pass the nonpositive true-upper check. Cell--cases are deterministic repeated diagnostics, not independent economic observations.}
\end{table}


The corrected gate and the deliberately false-null gate are both replayed. Every selected corrected action is feasible and has a nonpositive true-advantage upper endpoint. For the false-null policy, a strictly positive true-advantage lower endpoint identifies a harmful decision, not merely an uncertain one. The full catalogue retains accepted, blocked and unchanged actions, including cases where large non-nullity makes the verified gate conservative. These are deterministic diagnostics for known perturbed primitives. They are not additional cost-inference samples, an automatic discovery of learned null structure, or a guarantee under an unknown misspecified transition law.
